
\begin{abstract}
Long-term memory is essential for LLM agents that operate across multiple sessions, yet existing memory systems treat retrieval infrastructure as fixed: stored content evolves while scoring functions, fusion strategies, and answer-generation policies remain frozen at deployment. We argue that truly adaptive memory requires co-evolution at two levels: the stored knowledge and the retrieval mechanism that queries it. We present \ours{}, a self-evolving memory architecture that exposes its full retrieval configuration as a structured action space optimized by an LLM-powered diagnosis module. In each evolution round, the module reads per-question failure logs, identifies root causes, and proposes targeted configuration adjustments; a guarded meta-analyzer applies them with automatic revert-on-regression and explore-on-stagnation safeguards. This closed-loop self-evolution realizes an AutoResearch process: the system autonomously conducts iterative research cycles on its own architecture, replacing manual configuration tuning. Starting from a minimal baseline, the process converges autonomously, discovering effective retrieval strategies including entirely new configuration dimensions not present in the original action space. On LoCoMo, \ours{} outperforms the strongest baseline by 25.7\% relative and achieves a 78.0\% relative improvement over the minimal baseline. On MemBench, \ours{} exceeds the strongest baseline by 18.9\% relative. Evolved configurations transfer across benchmarks with positive rather than catastrophic transfer, indicating that the self-evolution process captures universal retrieval principles rather than benchmark-specific heuristics. Code is available at \url{https://github.com/aiming-lab/SimpleMem}.
\end{abstract}

